% P10-Evaluationsfeinschliff, 16.09.2026: Aussagegrenzen und Anforderungsbilanz; keine neuen Messwerte.
% P10-Core-Inhaltsprüfung, 15.09.2026: datierter Ergebnisanschluss; B-60 korrigiert.
% Belege: thesis-evidence/w2/core-content-audit-20260915/evaluation/
% B-59, 15.09.2026: datierte Ledger-Nachauswertung, Stützung und Reichweite.
% Beleg: thesis-evidence/w2/ledger-repetitions-20260915/ (Originale unverändert).
% P10-4a, 14.09.2026: Mechanismus, Quellenreichweite und v4-Nachweis angeglichen.
% M19 (12.09.2026): Leserführung und präzise Ergebnisdarstellung;
% Messwerte und historische Nachweisgrenzen bleiben erhalten.
% ============================================================
% §7.8 Gültigkeitsgrenzen — v01 (Claude, 07.09.)
% NEUE 8er-Struktur (Umsetzungsplan §7.8: knapp und konkret;
% beobachtete stille Lücken/schwache Hinweise/offene Pfade bleiben
% ERGEBNISSE in 7.3–7.5 und werden hier nicht erneut „entsorgt").
% Budget ~0,7 S., keine Tabelle.
% ============================================================
% HERKUNFT: Grenzen-Katalog = Umsetzungsplan §7.8 + Messprotokoll
% §2/§7 (Messgrenzen, Abgrenzung) + §6b (dokumentierte
% Kennzahl-Abweichung) + die in tex-v2 §§7.2–7.7 je Abschnitt
% benannten Einzelgrenzen (hier gebündelt, nicht dupliziert).
% ------------------------------------------------------------

\section{Gültigkeitsgrenzen}
\label{sec:eval-grenzen}

Die Aussagekraft der Evaluation hängt von ihren Quellen, den
Bewertungsverfahren und den untersuchten Systemständen ab. Die
folgenden Grenzen betreffen den Untersuchungsumfang; konkrete
Systemschwächen bleiben bei den jeweiligen Ergebnissen ausgewiesen.

\textbf{Quellen und Fälle.} Beide Quellenfälle stammen aus einer
Domäne, waren entwicklungsbekannt und lagen deterministisch
vorsegmentiert vor. Je Kombination aus Quellenfall und Konfiguration
wurde genau ein Hauptlauf bewertet. Zwei ergänzend ausgewertete
Ledger-Ausgaben erweitern nur den Stressfall; die drei Ausgaben
tragen eine beschreibende Streuungsangabe, keine belastbare Schätzung
allgemeiner Zuverlässigkeit. Für F fehlen entsprechende
Wiederholungen; die Streuung der Unterschiede zwischen den
Konfigurationen ist damit nicht bestimmt. Die fachliche Fallserie ist gezielt
gewählt und nicht repräsentativ. Auch die acht regelgebunden
ausgewählten Core-Fälle teilen teilweise Projektstände und umfassen
Testeingaben. Die Auswahl aus vorhandenen Vorschlägen erfasst nicht,
welche Anforderungen gar nicht erkannt wurden. Das Material trägt keine
Generalisierung auf fremde Quellen, Domänen oder Modelle; die
Modellreihe im Anhang bleibt eine ergänzende, nicht verifizierte
Beobachtung.

\textbf{Referenz und Urteil.} Erstellung und Bewertung des
Referenzbestands nutzen teilweise dieselben Werkzeuge und Beteiligten.
KI-Systeme lieferten den Referenzentwurf und die Inhaltsurteile;
der Autor konsolidierte den Bestand und war zugleich Entwickler und
Prüfer. Zusätzliche Kontrollen bestanden aus KI-Nachprüfungen
definierten Umfangs. Für die ergänzenden Auswertungen wurde mit Codex
ein Werkzeug desselben Anbieters wie das Systemmodell GPT-5.4 genutzt.
Getrennte KI-Prüfrunden gewährleisten keine voneinander unabhängigen
Urteile. Eine verblindete oder unabhängige menschliche
Zweitannotation fehlt. Die Korrektheitskennzahl weicht dokumentiert
von der ursprünglich protokollierten Gold-Precision ab
(Abschnitt~\ref{sec:eval-design}). Auch die konsolidierende Nachprüfung
war retrospektiv und kannte die bisherigen Urteile. Einzelne Matches
hängen von der Regelauslegung ab; die Auswirkungen alternativer
Auslegungen sind im Anhang beziffert. Die qualitative Lückentypisierung
besitzt keine validierte Wichtigkeitsskala. Bei der ergänzenden
Ledger-Auswertung wurden die Abdeckungsurteile des Hauptlaufs
übernommen und die neuen Urteile gegen dokumentierte Auslegungen
abgeglichen. Laufunterschiede und Bewertungsabhängigkeit lassen sich
daher nicht vollständig trennen. Die drei Autorenprüfungen betreffen
gezielt ausgewählte Grenzfälle dieser Ledger-Auswertung, keine
repräsentative Stichprobe. Die ergänzende Core-Inhaltsprüfung erfolgte
ebenfalls KI-gestützt. Ihre drei Autorenrückfragen ließen Auslegung
beziehungsweise Entstehungshintergrund offen; sie sind keine
menschliche Vollprüfung der Fälle.

\textbf{Queue und menschliche Bearbeitung.} Die
Adjudikations-Queue des Stresslaufs ist eine dokumentierte
deterministische Rekonstruktion; Aussagen über ihre historische
Anzeige und die tatsächliche Wirkung der Bedienaktionen tragen
nur die dafür vorhandenen Belege. Ein tatsächlicher
HITL-Durchlauf mit anschließender Neubewertung wurde nicht
ausgeführt --- das ist eine Untersuchungsgrenze, keine
Systemaussage; die Szenariorechnung in Abschnitt~\ref{sec:eval-luecken} bleibt eine bedingte
Rechnung. Davon zu unterscheiden sind \emph{Systembefunde} wie
fehlende Inhalte, die beruhigende Fehleinstufung ausgefilterter
Hinweise oder nicht auflösbare Herkunftspfade: Sie bleiben
Ergebnisse.

\textbf{Stände und Entstehung.} Die Läufe stammen aus historischen
Systemständen; ein finaler Integrationslauf fehlt. Ergänzende
Untersuchungen wie Stufenanalyse, Fallserie und Kontrolltests wurden
nach der Hauptauswertung vorgenommen. Ihre Regeln wurden vor der
jeweiligen Prüfung dokumentiert. Die Nachreview-Runde enthält dagegen
auch nachträgliche Auslegungsentscheidungen mit Kenntnis der Ergebnisse;
diese sind als solche datiert. Ebenso ist die spätere Überarbeitung
des Evaluationskapitels offengelegt. Das Herkunftsaudit~v4 entstand
nach der Diagnose unzureichend erfasster Herkunftswege. Seine
Endpunkte und nach Probeläufen präzisierten Regeln sind gesondert
dokumentiert (Anhang~\ref{anh:provenienz-v4}). Die Reproduktion
belegt die Wiederholbarkeit dieser Auswertung, keine Stabilität
der Modellausgaben. Ein Beitragspfad erklärt zudem nicht zwingend
alle Inhalte des erreichten Items; aktuelle Referenz- und
Änderungsprüfungen werden deshalb getrennt ausgewiesen.

Die Facettenvervollständigung
nachträglich aufgenommener Evidenz-Claims wurde erst nach Abschluss
der Messkampagne in den Gesamtablauf eingebunden
(Abschnitt~\ref{sec:evidenzgebundene-verarbeitung}); die hier ausgewerteten
Läufe entstanden vor dieser Einbindung. Auch die am 11.09.2026
nachträglich korrigierte Sprechersegmentierung und Fehlerausgabe
der Baseline-Stufe wurden gesondert technisch geprüft
(Anhang~\ref{anh:evaluation}); diese Nachweise erweitern die
Aussagekraft der historischen Messungen nicht.

Gemessen ist der maschinelle Bestand vor dem
Human-Gate beziehungsweise der jeweils bezeichnete Pfad. Die
integrierte Verarbeitung mit menschlichen Freigaben wurde anhand
ausgewählter historischer Fälle nachvollzogen. Einzelne Ablehnungen
verhinderten die Übernahme unvollständiger Vorschläge; der Einfluss
menschlicher Eingriffe auf die Ergebnisqualität wurde jedoch nicht
vergleichend untersucht. Nicht gemessen wurden der erforderliche
Prüfaufwand und die Leistungsfähigkeit eines einheitlichen finalen
Systemstands. Ungeprüft bleiben auch Betriebs- und Fehlerfälle
jenseits der geprüften Klassen
(unter anderem G01, G04, allgemeine Drift-Überwachung,
Absturz-Doppelwirkungen).

Die in den Ergebnisabschnitten berichteten stillen Lücken, schwach
relevanten Hinweise und offenen Provenienzpfade sind von diesen
Gültigkeitsgrenzen zu unterscheiden: Sie sind Ergebnisse der
Evaluation und bleiben als solche bestehen.
